\section{A numerical simulation and the individual upper bound}
\label{sec:upper}

Persistent symbolic marking must not become a numerical identity step.
We simulate each reuse by decreasing a finite label. The argument follows
the computability architecture of numbered higher-order path orderings
\cite[Section~5.3]{kop}, with that additional decrease made explicit.
It proves well-foundedness independently of the composition theorem.

\subsection{One finite arithmetic universe}
Fix a finite subtype-closed set $S$ and the ranked signature $\mathcal F$.
The auxiliary syntax has variables, lambdas, bottom $\bot_\tau$, ordinary
heads, numbered heads $f^k_\tau(\vec b;\vec e)$ for $k\in\N$, and a
separate ordinary auxiliary application $QU$. An ordinary head has label
$\omega$. Auxiliary application is distinct from the literal original
binary head $@_{\sigma,\tau}$.

Every node and binder type lies in $S$. In addition, every head's
\emph{original} result type and all original base-argument types lie in
$S$, even when a numbered head has been retyped. These restrictions give
finitely many possible native and application heads. Ranked heads belong
to the fixed finite signature. An ordinary term is in this universe
whenever its literal support is contained in $S$.

Use de Bruijn bound syntax and globally typed free names $(\tau,i)$,
where $\tau\in S$ and $i\in\N$. All finite contexts and their extensions
are represented in this one arithmetic carrier. Fresh opening chooses
a name outside a finite set. There is no application of a theorem about
one fixed context to unrepresented later context extensions.

Write $P^k_\rho=f^k_\rho(\vec b;\vec e)$. The raw auxiliary reduction
$\leadsto$ has the following root rules, always with the displayed
types and head-window conditions.
\begin{enumerate}
\item An ordinary $f(\vec b)$ reduces to $f^k(\vec b;\varnothing)$
for any finite $k$.
\item $P^{k+1}_\tau\leadsto P^k_\tau$.
\item If a base or extra argument $q$ has type
$\rho_1\to\cdots\to\rho_n\to\tau$, then
\[
 P^{k+1}_\tau\leadsto qP^k_{\rho_1}\cdots P^k_{\rho_n}.
\]
The right side uses ordinary auxiliary application.
\item If $f>g$ and $g$ belongs to the fixed head window, then
$P^{k+1}_{\operatorname{res}(g)}$ reduces to the ordinary $g$-head
whose base arguments are the appropriately retyped copies $P^k$.
\item If $f\approx g$, $g$ lies in the same window, and the base list
of $g(\vec c)$ is a primitive strict replacement of $\vec b$ as defined
below, then $P^{k+1}_{\operatorname{res}(g)}\leadsto g(\vec c)$.
\item $P^{k+1}_{\sigma\to\tau}$ reduces to
$\lambda x:\sigma.f^k_\tau(\vec b;\vec e,x)$ with fresh opening.
\item $(\lambda x.Q)U\leadsto Q[x:=U]$, and every nonbottom $Q:\tau$
reduces to $\bot_\tau$.
\end{enumerate}
Close these rules under lambda bodies and under each base and extra
argument of an ordinary or numbered head. Auxiliary application has
\emph{function-side closure only}: $Q\leadsto Q'$ gives
$QU\leadsto Q'U$. It has no argument-side closure.

For primitive replacement, a \emph{mark decrease} passes through a
common initial lambda spine and either changes an ordinary head's label
$\omega$ to a finite label, or changes a finite leading label $k$ to
$l<k$. It leaves the head, arguments, extras and output type fixed.
Variables, bottoms and auxiliary applications cannot be marked this way.
Every mark decrease is a nonempty raw path. A primitive multiset
replacement retains identical occurrences, removes a nonempty source
part, and inserts finitely many mark decreases of removed occurrences.
Deletion is allowed. In particular label zero cannot be marked again.

Quotient by the congruence that replaces equivalent heads and permutes
their permitted base arguments, preserving types, labels, lambda shape
and extra order. Its classes are finite and effectively enumerable in
the fixed head window. We may use least-code representatives. Let
$R_S$ be the positive-path relation induced by $\leadsto$ on these
classes. Both relations are arithmetic: finite raw paths and finite
equivalence witnesses are numerical certificates.

\subsection{Accessibility in one positive inductive definition}
\begin{lemma}[Arithmetic accessibility]\label{lem:arithmetic-accessibility}
For each fixed $S,\mathcal F$, the theory $\ID$ proves accessibility
of every well-typed auxiliary term under $R_S$.
\end{lemma}
\begin{proof}
Use two positive arithmetic operators, with typedness included:
\[
 I(Q)\ \longleftrightarrow\
   \forall Q'\,(Q R_S Q'\Rightarrow I(Q')),
 \qquad
 J(M)\ \longleftrightarrow\
   \forall M'\,(M\DM(R_S)M'\Rightarrow J(M')).
 \tag{A}
\]
Here $\DM(R_S)$ is ordinary strict finite multiset replacement, retaining
identical occurrences and replacing removed ones by $R_S$ descendants.
Neither clause refers to the other's inductive predicate. Disjoint
numerical tags combine them into one positive arithmetic definition.
The $I$-stages of that definition equal the standalone accessibility
stages, since its clause never consults the $J$ tag.

For each of the finitely many types in $S$ define, by external recursion
on type,
\[
 C_o(Q)=I(Q),\qquad
 C_{\sigma\to\tau}(Q)\iff
       \forall U\,(C_\sigma(U)\Rightarrow C_\tau(QU)).
 \tag{T}
\]
Each is a first-order formula in the language with $I$; no further
inductive definition is introduced at an arrow. Quantification over
computable typed arguments abbreviates a finite disjunction over $S$.

We give the computability facts used in the constructor induction.
Joint induction on type proves that a non-lambda is computable whenever
all its direct reducts are computable; that variables and bottoms are
computable; that computability is preserved by $\leadsto$ and $R_S$;
and that computability implies $I$. At base type, closure of (A) and
transitivity of positive paths give the assertions. At an arrow, $QU$
has no top beta step when $Q$ is not a lambda. Apart from bottom its
reducts reduce the function $Q$, so codomain induction applies. Bottom
and variable applications are handled by this same neutral argument.
Function-side closure gives downward closure of (T).

For the implication from computability to $I$ at an arrow, apply the
term to a fresh computable variable $x$. Every positive path
$Q R_S Q'$ gives $Qx R_S Q'x$. Induction on $I(Qx)$ therefore proves
$I(Q)$: all descendants $Q'$ inherit the corresponding strictly
smaller accessibility premise.

The lambda fact is also needed. Suppose every computable instance
$V[x:=U]$ is computable. Substituting a fresh variable first gives
$I(V)$. Induct on this accessibility while retaining the universal
instance premise. The direct reducts of $(\lambda x.V)U$ are its given
beta instance, a function-side body reduct, or bottom. For a body
reduct, substitution takes a raw step to a raw path or equality; the
only possible loss of a step is harmless, as when a variable-to-bottom
step becomes bottom-to-bottom. Mark decreases commute with substitution
because their leading labels remain unchanged. Downward closure and the
induction hypothesis therefore make every reduct computable. The
neutral argument gives the application and hence the lambda.
All statements respect the finite equivalence quotient.

Next prove that accessibility of finitely many members implies $J$ of
their multiset. First prove closure of $J$ under multiset union by outer
$J$-induction on one component and inner $J$-induction on the other.
Label removed occurrences by their component of origin and assign each
new occurrence to a removed dominator. The target splits into two
components, each unchanged or decreased, with at least one decreased.
The nested induction and downward closure supply its $J$-membership.
Now $I$-induction on $Q$ gives $J(\{Q\})$: any replacement consists of
finitely many descendants $Q_i$, the induction hypothesis gives each
$J(\{Q_i\})$, and finite union gives the replacement's membership.
Another finite union proves the assertion for arbitrary accessible
member multisets.

Finally prove that a constructor with computable arguments is computable
by the lexicographic induction
\[
 (\text{head class},\ \text{base multiset},\
    \text{label in }\N\cup\{\omega\},\ \text{extra multiset}).
 \tag{L}
\]
Use finite induction on head classes, $J$-induction on the base list,
ordinary induction on finite labels followed by the $\omega$ case,
and $J$-induction on extras. Every earlier invariant quantifies all
later parameters, including the output type and the finite typed extra
list. These are first-order formulas in the language of $\ID$.

The neutral lemma reduces this assertion to its direct reducts.
Ordinary marking decreases $\omega$ to a finite label, and the added
lowering rule strictly decreases a finite label. Select applies a
computable selected argument to computable lower-label self-copies.
Copy first obtains those copies and then uses a smaller head class;
its target label may reset to $\omega$ because the first component
has decreased. In the multiset rule, every inserted argument is a
positive mark-decrease descendant of a removed one. Downward closure
makes it computable, and the base multiset decreases under $\DM(R_S)$.
This earlier decrease again permits the outer label to reset to
$\omega$. Abstraction creation decreases the label before adding a
fresh extra, and the lambda fact applies to every computable
substituted extra. A base or extra argument reduction decreases its
corresponding multiset. Bottom is computable. These exhaust the rules.

Structural induction under a computable substitution now makes every
auxiliary term computable and hence accessible. All syntax, typing,
paths, multiset assignments and invariants used above are arithmetic
in the one fixed universe, so this is an $\ID$ proof of the stated
individual memberships.
\end{proof}

\subsection{Finite budgets and preservation of the universe}
A symbolic marked term is $S$-admissible if its node and binder types
and every head's original declaration lie in $S$. Ordinary endpoints
are admissible for $S=\Ty(s)\cup\Ty(t)$.

\begin{lemma}[Locality]\label{lem:window-locality}
Every term in a finite weak or strict derivation with $S$-admissible
endpoints, and every numerical intermediate constructed below, is
$S$-admissible.
\end{lemma}
\begin{proof}
Induct over the actual comparison and occurrence-list derivations.
Var and matching introduce no type. Abs opens at the given binder
type. Put and Strict change only the mark. FAbs appends the domain
of the existing outer arrow and uses its codomain.

For Select, each supplied domain and every successive arrow suffix is
a subtype of the selected argument's type. A self-copy retains its
original head declaration. The full application consumes original
lambdas, substitutes admissible copies, and possibly marks a residual
head and appends those copies. Induction on that literal operation
introduces no new type or head. This establishes admissibility of the
continuation source before its derivation induction is applied.

Copy and Stat retype only to actual target-argument types. Their new
target head is present in the endpoint and therefore satisfies the
head-window guard. Their occurrence witnesses use actual endpoint
arguments. This proves the symbolic assertion.

Numerical realization retains this typed shape. Lowering and mark
decrease change no declaration. Copy and multiset reset use the
already admissible target head. Auxiliary Select applications and the
lambdas created by the full-tuple helper have the domain and suffix
types already checked above. Beta substitution, contexts and equivalent
head permutations preserve the same window. Thus every constructed
raw edge belongs to this single $R_S$ universe.
\end{proof}

Write $q\sim_n Q$ when $Q$ realizes the symbolic source $q$ with the
same typed shape, ordinary heads at label $\omega$, and every marked
head at some finite label at least $n$. The realization preserves base
and extra lists, commutes with simultaneous substitution, and weakens
to every smaller threshold. For an ordinary $q$ there is a unique
$\tr(q)$, available at every threshold.

\begin{lemma}[Full-tuple simulation]\label{lem:tuple-simulation}
If $\ap q{a_1,\ldots,a_l}$ is defined, sufficiently high realizations
of $q,a_i$ reduce by a finite auxiliary path to a realization of that
exact result at any prescribed lower threshold $n$. The helper spends
at most $l+1$ units beyond that threshold.
\end{lemma}
\begin{proof}
Consume the same original lambda spine by beta steps, with the supplied
arguments substituted in their common environment. If the tuple ends,
this is the required result. Otherwise source definedness says that
the original residual head is functional. Mark an ordinary head to
a sufficiently large finite label, or retain its available finite
label. For each remaining argument, create an abstraction and
beta-contract it, thereby appending that argument while decreasing
the label by one. There are at most $l$ such decreases. Substituted
old arguments retain their own thresholds. A protected-variable
overapplication is undefined in the source and is never a premise.
All helper types remain in the window by Lemma~\ref{lem:window-locality}.
\end{proof}

\begin{lemma}[Positive simulation]\label{lem:positive-simulation}
For every finite weak derivation $\W q t$ in a fixed window there is
a computable finite budget $b$ such that every $q\sim_b Q$ reduces,
possibly by an empty path and modulo finite equivalence, to $\tr(t)$.
For a strict derivation the budget can be chosen so that the path is
positive. Consequently, for ordinary confined endpoints,
\[
 s\cs t\quad\Longrightarrow\quad
 \tr(s)R_{\Ty(s)}\tr(t).
\]
\end{lemma}
\begin{proof}
All original targets are ordinary by Lemma~\ref{lem:elementary}.
Induct jointly on the finite comparison and list witnesses, taking
maxima of their finitely many required budgets.

Var requires no step. Abs lifts a body path; Fun performs child paths
inside base arguments and uses the finite head/permutation equivalence.
Put chooses a finite leading label above its child's budget.
Select reserves one unit for the Select rule, the full-tuple cost
of Lemma~\ref{lem:tuple-simulation}, and the continuation budget.
FAbs reserves one unit and lifts its body path. Copy reserves one
unit, builds its ordinary target from lower-label self-copies, and
performs the finite child paths in its arguments.

For Stat retain the specified source occurrences unchanged in the
intermediate target. Each inserted occurrence has a strict premise,
hence a defined mark and a shorter weak derivation from that mark.
If its numerical source argument is ordinary, choose an
$\omega$-to-finite mark above the child budget. If it is already
numbered, reserve an additional unit and strictly lower the finite
label. Under initial lambdas use the same context step. These are
exactly the permitted primitive replacements. The numerical multiset
rule builds an ordinary target whose retained arguments then follow
their weak paths and whose inserted arguments follow their marked
weak paths.

For Strict itself perform this same nonempty mark decrease on its
source before following its weak premise. At a symbolically
already-marked source the symbol stays the same but its finite
numerical label strictly decreases. This makes the simulated path positive. Label zero
is never reused because the budget reserves the needed unit.

All choices depend only on the finite proof certificate. Locality
keeps every intermediate in the original endpoint window. For
ordinary $s$, $\tr(s)$ has only $\omega$ labels and meets every
finite budget. Confinement makes that window precisely $\Ty(s)$.
The marked-target auxiliary relation of Appendix~\ref{app:composition}
is not part of this simulation; its reflexive marked strict judgments
are never interpreted as positive paths.
\end{proof}

We use the following published ordinal-analysis input. For a positive
arithmetic inductive operator, every individual membership atom provable
in $\ID$ has inductive norm strictly below the Bachmann--Howard ordinal.
With stages beginning at the operator applied to the empty set, the
norm for the accessibility operator is exactly descending rank
\cite[Definitions~6.3.1, 6.3.5 and 9.1.2,
Lemma~6.5.2, Theorems~9.5.5--9.5.6]{pohlers}.

\begin{theorem}[Well-foundedness and individual bound]\label{thm:upper}
The confined strict relation is well founded, and
$\rk_{\cs}(s)<\BH$ for every ordinary term $s$.
\end{theorem}
\begin{proof}
Every descendant of $s$ remains in $S=\Ty(s)$.
Lemmas~\ref{lem:arithmetic-accessibility} and
\ref{lem:positive-simulation} exclude an infinite strict chain.
Their positive-edge simulation gives
\[
 \rk_{\cs}(s)\le \rk_{R_S}(\tr(s)).
\]
The arithmetic accessibility proof supplies an $\ID$ proof of the
individual $I$-membership of $\tr(s)$ in this fixed window. The tagged
$I$ component has its standalone accessibility norm. Applying the
published norm theorem yields
$\rk_{R_S}(\tr(s))<\BH$, as required. No transitivity assumption
about the symbolic comparison is needed in this argument.
\end{proof}

\section{Native terms with cofinal lower ranks}\label{sec:lower}

For a closed native term $t:A$, let $r_A(t)$ be its rank in the
restriction of $\cs$ to closed native terms of that same type.
These ranks are countable by well-foundedness and countability of
syntax. Restriction cannot increase rank. The lower construction
uses these ranks themselves, without a rank borrowed from the
auxiliary numerical system.

\subsection{A carrier inside literal T}
Fix closed native terms $P,K:B\to B$ and $U:B$ with $\W K P$
and $\Ty(P)\subseteq\Ty(K)$. Put
$H=o\to B\to B$, $F=\lambda n:o.K$ and
\[
 W=R^B0FU,\qquad U_1=R^B0,\quad U_2=U_1F.
\]
For any required result type $\rho$ define proof intermediates
\[
\begin{aligned}
 Q_\rho&=@^\star_\rho(U_2,U),\\
 C_{2,\rho}&=@^\star_\rho(U_1,F;Q_B),\\
 C_{1,\rho}&=@^\star_\rho(R^B,0;C_{2,H},C_{2,B}),\\
 C_{0,\rho}&=(R^B)^\star_\rho(;C_{1,o},C_{1,H},C_{1,B}).
\end{aligned}
\]
They are marked comparison terms, not additional native constructors.

\begin{lemma}[Native carrier]\label{lem:carrier}
$W\cs U$, and $W\cs T$ implies $W\cs PT$ for every ordinary
closed $T:B$ carrying the stated support guard.
\end{lemma}
\begin{proof}
Select the second base argument of $Q_B=\mk(W)$ to obtain $U$.
For the second assertion, the hypothesis is a weak proof from $Q_B$
to $T$. Select the extras $C_{1,B},C_{2,B},Q_B$ in turn to prefix
that proof by a derivation from $C_{0,B}$.

To derive the other Copy premise, start from $C_{0,B\to B}$.
Select its extra $C_{1,H}$ with one self-copy of type $o$; this
appends that copy to the marked $C_1$ owner. Select its original
$C_{2,H}$ extra with one further self-copy, and then its original
base argument $F=\lambda n.K$ with one self-copy. The result is
literally $K$, since $n$ is unused. Append the given weak proof
from $K$ to $P$. Copy at $R^B>@$ now constructs $PT$ from these
two whole-source premises.

Three outer Select prefixes carry this proof back to $Q_B$.
They select $U_2,U_1,R^B$ with respectively one, two and three
simultaneous copies at their displayed types. Strict marking gives
$W\cs PT$. The support of $T$ is contained in that of $W$ by
hypothesis; the support of $P$ occurs through $K$, and the final
application adds only $B,B\to B$, already present. Thus every
comparison is confined. No composition theorem is used.
\end{proof}

Let the semantic base type be $\iota$, interpreted as the countable
ordinals, and translate it to the actual type $\sigma=o\to o$.
Translate arrows homomorphically. Define erasable templates
\[
 E_\iota(i)=i,\qquad E_{A\to D}(i)=\lambda x:|A|.E_D(i).
\]
For $B=|A|$, $\eta=B\to B$, and fixed $P:\eta,U:B$, put
\[
 K_i=((\lambda a:\eta.\lambda b:\eta.a)P)E_{A\to A}(i),
 \qquad
 W_i=R^B0(\lambda n:o.K_i)U.
\]

\begin{lemma}[Indexed carrier]\label{lem:indexed-carrier}
$W_i\cs U$, and $i\cs j$ implies both
$W_i\cs W_j$ and $W_i\cs PW_j$.
\end{lemma}
\begin{proof}
First $K_i\cs P$ has a direct proof. Mark the outer application and
select its first base argument with one self-copy. The result is the
marked inner application retaining $P$ as its second base argument.
Select $P$ with an empty tuple. Since $P$ is a subterm, confinement
also holds. Strict inclusion supplies the weak hypothesis of
Lemma~\ref{lem:carrier}.

The context defining $W_i$ contains $i$ in its erasable argument.
Context stability therefore gives $W_i\cs W_j$ from $i\cs j$.
Apply the carrier lemma to that target to get $W_i\cs PW_j$.
The common context contributes the same types; the support of $j$
is contained in that of $i$. The final application uses only types
already present. The selector proof did not compose beta reductions.
\end{proof}

\subsection{Rank realization and full-prefix iteration}
Semantic arrow types contain all set functions of the indicated types,
not only monotone ones. Countable suprema are pointwise at arrows.
For closed actual terms define
\[
 t\models_\iota\alpha\iff r_\sigma(t)\ge\alpha,
\]
and let $t\models_{A\to D}f$ mean that for every semantic $x:A$
and every closed $u\models_A x$, the ordinary application $tu$
realizes $f(x)$.

\begin{lemma}[Realization closure]\label{lem:realization}
Realization is preserved when a term is replaced by a strict predecessor.
It is downward closed in its semantic value, with pointwise order at
arrows. If one fixed term realizes every member of a nonempty countable
family, it realizes their pointwise supremum. The typed lambda
fundamental lemma holds for these realizers.
\end{lemma}
\begin{proof}
Induct on the semantic type. At $\iota$, use the rank inequality
for a strict edge and the defining property of ordinal supremum.
At an arrow, apply the induction hypotheses to each realizing
argument; context stability gives $tv\cs uv$ from $t\cs u$.
The other assertions are pointwise. For the lambda lemma, its
application beta-reduces strictly to the substituted body, so strict
predecessor closure transports the body's realization backward.
An induction on the number of beta steps permits any finite sequence.
\end{proof}

The closed native iterator is
\[
 \It_A=\lambda i:\sigma.\lambda p:\eta.\lambda u:B.\
 R^B0\bigl(\lambda n:o.((\lambda a:\eta.\lambda b:\eta.a)p)
                   E_{A\to A}(i)\bigr)u.
\]
For an arbitrary semantic $f:A\to A$ define
\[
 f^0(x)=x,\qquad
 f^{\beta+1}(x)=f(f^\beta(x)),\qquad
 f^\lambda(x)=\sup_{\xi<\lambda} f^\xi(x)
 \quad(\lambda>0\text{ limit}).
 \tag{I}
\]

\begin{theorem}[Native finite-type iteration]\label{thm:iteration}
$\It_A$ realizes $(\alpha,f,x)\mapsto f^\alpha(x)$ with convention
\textup{(I)}, even when $f$ is not monotone.
\end{theorem}
\begin{proof}
Fix closed realizers $P\models f$ and $U\models x$. For each closed
index $i:\sigma$, prove that the single term $W_i$ realizes every
orbit value $x_\beta=f^\beta(x)$ for $\beta\le r_\sigma(i)$.
Use outer well-founded induction on $r_\sigma(i)$ and inner induction
on $\beta$.

At zero, $W_i\cs U$ supplies the conclusion. If
$\beta=\gamma+1\le r_\sigma(i)$, the rank equation gives a closed
native $j$ with $i\cs j$ and $r_\sigma(j)\ge\gamma$.
Otherwise all successor ranks would be at most $\gamma$, contrary
to $r_\sigma(i)\ge\gamma+1$. The outer induction makes $W_j$
realize $x_\gamma$. Hence $PW_j$ realizes $f(x_\gamma)$, and
$W_i\cs PW_j$ transports that realization to $W_i$.

At a nonzero limit the same $W_i$ realizes every earlier $x_\xi$.
There are countably many earlier stages. The supremum clause of
Lemma~\ref{lem:realization} gives $x_\beta$. This uses the entire
prefix; a cofinal subsequence would not suffice for an arbitrary
nonmonotone orbit.

If $i\models\alpha$ then $\alpha\le r_\sigma(i)$.
Finally transport the realization backward through the finitely
many beta steps from $\It_A iPU$ to this carrier. All ranks and
edges are those of the closed native confined relation.
\end{proof}

\subsection{Explicit ordinal operations and published calibration}
Let $Z=\lambda z:o.z:\sigma$ and define
\[
 s=\lambda a:\sigma.((\lambda u:\sigma.\lambda v:\sigma.u)Z)a.
\]
$Z$ realizes zero. The body of $sa$ lies strictly above $a$ by
outer-argument Select, so beta transport makes $s$ realize right
successor. Put
\[
 \operatorname{Cap}(P,U)=R^B0(\lambda n:o.P)U.
\]
The same carrier is strictly above every finite $P^kU$, by induction
using Lemma~\ref{lem:carrier}. Consequently
$w=\operatorname{Cap}(s,Z)$ realizes $\omega$.

Writing $\It=\It_\iota$, define the following closed abbreviations:
\[
\begin{aligned}
 \mathrm{add}&=\lambda a\,b.\It\,b\,s\,a,\\
 \mathrm{mul}&=\lambda a\,b.\It\,b\,(\lambda x.\mathrm{add}\,x\,a)\,Z,\\
 \mathrm{one}&=sZ,\\
 \mathrm{pow}&=\lambda a.\It\,a\,(\lambda x.\mathrm{mul}\,x\,w)\,\mathrm{one},\\
 \mathrm{fix}&=\lambda f\,z.\It\,w\,f\,(s z).
\end{aligned}
\]
Their variables have the unique indicated ordinal-realizer types.
They realize, respectively, $\alpha+\beta$, $\alpha\beta$,
$1$, $\omega^\alpha$, and $f^\omega(\zeta+1)$.
The formulas specify the order of the noncommutative operations:
multiplication iterates addition on the right, and exponentiation
iterates right multiplication by $\omega$.
Let $\mathrm{Next}=\mathrm{fix}\,\mathrm{pow}$, realizing the
function that takes $\zeta$ to the least epsilon number above it.

Put $A_0=\iota$ and $A_{l+1}=A_l\to A_l$. The finite-type
extractors of Simmons \cite[Defs.~2.10--2.14]{simmons2004}
are
\[
 B_l h h_l\cdots h_1
   =\Fix\bigl(\lambda\alpha.\ h^\alpha h_l\cdots h_1\,0\bigr),
 \qquad B_l:A_{l+2}.
 \tag{B}
\]
They have actual closed native realizers. For $l=0$, the inner body
is $\It_\iota\,\alpha\,h\,Z$. For $l\ge1$, it is
\[
 (\It_{A_l}\,\alpha\,h\,h_l)h_{l-1}\cdots h_1Z.
 \tag{B$_T$}
\]
Prefix the corresponding typed lambdas and use $\mathrm{fix}$ for
$\Fix$. The terminal $Z$ is present in every case. The iterator
theorem and the lambda fundamental lemma verify the realization
of (B). All names here abbreviate finite expressions in the original
native constants, lambdas and applications.

We state the published calibration precisely. Let $\Omega$ be the first
uncountable ordinal, set $a_0=0$ and $a_{n+1}=\Omega^{a_n}$, and let
$\Lambda=\sup_n a_n=\varepsilon_{\Omega+1}$.
Simmons's finite-type iteration uses exactly the full-prefix pointwise
convention (I), on the full function spaces. His helpfulness theorem
for every $B_l$ is supplied in
\cite[Theorem~5.10]{simmons2008}, completing the hypotheses used
in \cite[Theorem~6.1 and Corollary~6.2]{simmons2004}.
For the Veblen base pair put
\[
 N_B=B_1B_0\mathrm{Next},\qquad \epsilon_B=\Gamma_0=N_B(\omega).
\]
The displayed identity following
\cite[Lemma~6.6]{simmons2008} gives the last equality.
The resulting sequence
\[
 \Delta_0=\Gamma_0,\quad
 \Delta_1=N_B\Gamma_0,\quad
 \Delta_{l+2}=B_lB_{l-1}\cdots B_0N_B\Gamma_0
 \tag{D}
\]
satisfies $\Delta_n=\psi_B(a_n)$ by the cited calibration.
Here $\psi_B$ is the Veblen-hull collapse, not the arithmetic-hull
collapse attached to $(\mathrm{Next},\varepsilon_0)$.

For completeness, the hypothesis called tameness in that calibration
holds for the displayed towers: their recursive coefficient norms are
at most one, whereas every value in the critical base sequence is at
least $\Gamma_0>1$. This is the actual inequality needed in the proof;
no assertion that a critical ordinal is finite is used.

Replace the operations in (D) by their native realizers and call the
resulting closed terms $t_n$. Each has actual outer type $\sigma=o\to o$,
and
\[
 r_\sigma(t_n)\ge\psi_B(a_n).
 \tag{R}
\]
The native syntax alone supplies these witnesses.

\subsection{The endpoint is the same Bachmann--Howard ordinal}
We record why the tower values are cofinal at the endpoint used for the
upper bound. In the Veblen convention, an $N_B$-critical ordinal is an
orbit value $N_B^\alpha(\Gamma_0)$, not a fixed point of the strictly
inflationary function $N_B$. These criticals are closed under countable
increasing suprema. Section~5 of \cite{simmons2004} characterizes
$\psi_B(\xi)$ as the least such critical $\theta$ for which
\[
 \eta<\xi,\quad |\eta|<\theta
       \quad\Longrightarrow\quad \psi_B(\eta)<\theta .
 \tag{P}
\]
The coefficient norm is recursively the maximum of the countable
coefficients in the base-$\Omega$ representation and of the norms
of its exponents. In particular $|a_n|\le1$.
The Veblen-hull endpoint $\psi_B(\varepsilon_{\Omega+1})$ is the
$\BH$ of the inductive-norm calibration
\cite[Sections~9.4--9.5]{pohlers}.

\begin{lemma}[Tower cofinality]\label{lem:tower-cofinality}
$\sup_n\psi_B(a_n)=\psi_B(\Lambda)=\BH$.
\end{lemma}
\begin{proof}
Let $\theta=\sup_n\psi_B(a_n)$. For $n<m$, the norm bound
$|a_n|\le1<\psi_B(a_m)$ and (P) give
$\psi_B(a_n)<\psi_B(a_m)$. The countable supremum is therefore
an $N_B$-critical ordinal.

If $\eta<\Lambda$ and $|\eta|<\theta$, choose $n$ large enough
that both $\eta<a_n$ and $|\eta|<\psi_B(a_n)$. Such a common
index exists by the two increasing cofinal sequences. Applying
(P) at $a_n$ gives $\psi_B(\eta)<\psi_B(a_n)\le\theta$.
Thus $\theta$ satisfies the endpoint clause, and its least
critical solution obeys $\psi_B(\Lambda)\le\theta$.
Conversely (P) at $\Lambda$ applies to each $a_n$, since
$|a_n|\le1<\psi_B(\Lambda)$, and gives
$\psi_B(a_n)<\psi_B(\Lambda)$. Take their supremum.
This uses the defining clause of this particular collapse; it
does not assume monotonicity for some different collapse convention.
\end{proof}

\begin{theorem}[Exact descending height]\label{thm:height}
The descending height of $\cs$ is $\BH$, already cofinally witnessed
by closed native T terms of outer type $o\to o$.
\end{theorem}
\begin{proof}
The realizers (R) and Lemma~\ref{lem:tower-cofinality} give the
lower bound for the restriction to closed native terms. Its ranks
are bounded by the ranks in the full ordering.
Theorem~\ref{thm:upper} bounds each individual full rank below
$\BH$. Taking the supremum of ranks plus one gives equality.
No individual term has been assigned rank $\BH$.
\end{proof}
